% diagrams/firstwords/papeleria.tex -- en la papelería: un cuaderno azul,
% lápices morados y un bolígrafo nuevo para Sam.
\begin{tikzpicture}[dibujofw]
  % el cuaderno, con su espiral
  \filldraw[fill=white] (-5.2,0) rectangle (-1.2,5.2);
  \foreach \y in {0.5,1.1,...,4.7} {\draw (-5.45,\y) arc[start angle=270, end angle=90, radius=0.18] -- ++(0.35,0);}
  \filldraw[fill=white] (-4.4,3.6) rectangle (-2.0,4.4);
  % los lápices
  \foreach \x/\a in {0.2/-8, 1.3/0, 2.4/8} {\begin{scope}[shift={(\x,0)}, rotate=\a]
    \filldraw[fill=white] (-0.3,0.6) rectangle (0.3,4.4);
    \filldraw[fill=white] (-0.3,4.4) -- (0,5.3) -- (0.3,4.4) -- cycle;
    \fill (0,5.3) -- (-0.1,5.0) -- (0.1,5.0) -- cycle;
    \filldraw[fill=white] (-0.3,0.2) rectangle (0.3,0.6);
  \end{scope}}
  % el bolígrafo de Sam
  \begin{scope}[shift={(4.4,0.2)}, rotate=-5]
    \filldraw[fill=white] (-0.28,0) rectangle (0.28,4.0);
    \filldraw[fill=white] (-0.28,4.0) -- (0,4.8) -- (0.28,4.0) -- cycle;
    \filldraw[fill=white] (-0.32,-0.9) rectangle (0.32,0.8);
    \filldraw[fill=white] (0.32,-0.5) rectangle (0.48,0.6);
  \end{scope}
\end{tikzpicture}
